\section{Scientific history of this result}\label{scientific-history-of-this-result}

This section records how the reported specifications and results came to be what they are. It is here so that the argument above does not have to carry it, and so that a reader comparing this version with an earlier one can see what changed and why. Nothing in it is a competing model.

\textbf{The consumption specification.} Earlier drafts reported a specification in which the consumption coefficient was fixed at one as a numeraire, and one in which the single-adult consumption curvature was estimated at \(\theta_c\approx0.168\). Neither is the reported model. The current specification fixes the shock scale, fixes the consumption curvature at exactly zero, and estimates the consumption weight instead. That change improves the criterion by 23.753 log-points for single adults and 15.659 for couples, and it changes the welfare aggregator from an arithmetic mean of reachable consumption to a power mean of order \(\beta_c\). Every welfare quantity computed under the earlier convention is superseded, and the corresponding figures were regenerated rather than relabelled.

\textbf{The sample.} The predecessor frames contained 1,555 single-adult and 2,275 couple households. Twelve single-adult and fifty-two couple households were removed when the hours support and the occupation mapping were applied to the observed rows, and three further single-adult households were removed because their own observed job prices to non-positive disposable consumption. The current descriptive tables are computed on the resulting 1,540 and 2,223 households, not on the predecessor frames, and the funnel in Section 2 ends where the estimation begins.

\textbf{The disposable-income convention.} Single-adult disposable income was previously aggregated over the decider only; it is now aggregated over all resident household members, which is the convention the couples application always used. The two applications are now on the same accounting convention.

\textbf{The decomposition result.} An earlier draft reported household endowments and needs as the largest component for single adults, and a one-factor figure of about seventy-seven per cent. Both are withdrawn. The corrected computation gives 49.16 per cent to job access against 34.20 to resources and needs, and the corresponding one-factor effect for all non-preference circumstances is 67.2 per cent.

\textbf{The bridge between references.} The signed gap between the flat-consumption reference and the non-work-bundle reference was previously reported with the wrong sign for couples. The premise audit in Section 6 locates the cause in the normalization of the opportunity kernel on the domain the reference mass uses, and the corrected bridge satisfies the inequality the premises require in both samples.

\textbf{The consumption normalizer.} Three values of \(\lambda_c\) were in circulation because different panels recomputed it over their own rows. Under exact log consumption the constant cancels from the money metric, so no result changed; the paper now reports one constant per population.

